\documentclass[10pt,a4paper]{amsart}
\usepackage[margin=19mm]{geometry}
\usepackage{amsmath,amssymb,mathtools}
\usepackage[scr=rsfs,cal=euler]{mathalfa}
\usepackage{xfrac}
\usepackage[unicode,hidelinks,bookmarks=false]{hyperref}
\newcommand{\ie}{i.e.\ }
\newcommand{\cf}{cf.\ }
\newcommand{\resp}{respectively\ }
\newcommand{\Subsection}{\subsection}
\providecommand{\nobreakdash}{\nobreak\hbox{-}}
\setlength{\emergencystretch}{2em}
\multlinegap0pt
\usepackage{bm}
\usepackage{bbm}
\usepackage{dsfont}
\usepackage{xspace}
\usepackage{cancel}
\usepackage{mathtools}
\usepackage{amsthm}
\makeatletter
\@ifundefined{lem}{\newtheorem{lem}{Lemma}}{}
\makeatother
\title[The Optimality of a Nested Generalized Pairwise Group Testin]{The Optimality of a Nested Generalized Pairwise Group Testing Procedure}
\author{Yaakov Malinovsky, Viktor Skorniakov}
\date{Source: arXiv:2506.15797v1 (CC BY 4.0)}
\begin{document}
\maketitle

\noindent\emph{Source and licence.} The Optimality of a Nested Generalized Pairwise Group Testing Procedure. Version arXiv:2506.15797v1 (CC BY 4.0), distributed under the Creative Commons Attribution 4.0 International licence, as recorded in the arXiv licence metadata for this exact version and shown on the abstract page https://arxiv.org/abs/2506.15797v1. Full rights notice: licenses/arxiv-2506.15797-CC-BY-4.0.txt.\par\smallskip

\noindent\textit{Editorial scope.} Counted unit: the Lem whose source label is \texttt{l:OAT\_structure} in the article (source0.tex lines 482--489, source label \texttt{l:OAT\_structure}), delivered together with the article's own complete original proof of it (source0.tex lines 491--530, one \texttt{proof} environment of 40 lines). No mathematics is written, repaired or re-derived: statement and proof are the article's own text line for line. Required context retained uncounted: l:single\_defective\_proc\_as\_OAT (lines 395--401). Every definition, display and label of those lines is retained unchanged. Omitted from this package and not redistributed: 1--394, 402--481, 490--490, 531--845. Nothing inside a retained excerpt is omitted or altered: every retained body is delivered exactly as the article's lines are written, so no omission receipt of any kind accompanies this package. Citations: the retained excerpts carry no citation command at all, so no bibliography entry of the article is transcribed and no reference is redistributed. Reference adaptation: none was needed and none was made; every reference inside the delivered unit resolves inside it, so no unresolved reference or citation is produced. Printed slips kept literal: no printed word, symbol, hypothesis or number of the counted statement or of its proof was altered. Layout and class: the article's own class is \texttt{article} with packages including amsmath, graphicx, psfrag, epsf, enumerate, natbib, amsthm, amsfonts, epsfig, amssymb, eucal, booktabs, multirow, siunitx; the wrapper is the lane's pinned \texttt{amsart} editorial wrapper at 10pt on A4, which restates the theorem-like environments the retained text needs and adds the article's own macro definitions that the retained text needs (none), with extra wrapper packages (bm, bbm, dsfont, xspace, cancel, mathtools). The delivered numbering is the unit's own. Assets: no retained span contains a figure environment, a graphics-inclusion command, a PDF-page inclusion, a table or a TikZ picture; no image file is copied into this package, the package declares no asset dependency and no image file is redistributed with it. Licence: version arXiv:2506.15797v1 of this article is distributed under the Creative Commons Attribution 4.0 International licence, as recorded in the arXiv licence metadata for this exact version and shown on the abstract page https://arxiv.org/abs/2506.15797v1. A full rights notice is delivered at licenses/arxiv-2506.15797-CC-BY-4.0.txt. Characters: every retained excerpt is pure ASCII, so the delivered engine, XeLaTeX with its default Latin Modern text font, reports no missing character.
\begin{lem}\label{l:single_defective_proc_as_OAT}
Assume that the initial group $(u_1, \dots, u_n)$ is contaminated. If, at each testing stage, a nested procedure is allowed to test only subsets consisting of units from the contaminated set, preserving their order from left to right, then the optimal procedure for isolating one defective unit is given by an OAT corresponding to the sequence of weights
\[
w_1 = \frac{p_1}{\alpha}, \quad w_2 = \frac{p_2 q_1}{\alpha}, \quad \dots, \quad w_n = \frac{p_n q_1 \cdots q_{n-1}}{\alpha}, \quad \text{where } \alpha = 1 - q_1 \cdots q_n.
\]
The isolated unit is the first defective encountered in the sequence $u_1, u_2, \dots, u_n$ when inspected from left to right.
\end{lem}

\begin{lem}\label{l:OAT_structure}
Consider the optimal procedure described in Lemma \ref{l:single_defective_proc_as_OAT}. Assume that
\begin{align}
\label{eq:con1}
1 - \frac{1}{\sqrt{2}} < p_1 \leq p_2 \leq \dots \leq p_n < \frac{3 - \sqrt{5}}{2}.
\end{align}
Then, for \( n \geq 3 \), in the first testing step, the procedure tests at most the first \( n - 2 \) units, and the penultimate unit is tested individually.
\end{lem}

\begin{proof}
Consider the last three terminal nodes of the OAT described in Lemma~\ref{l:OAT_structure}, which implement the optimal procedure. Their weights are given by
\[
w_{n-i} = \alpha^{-1} p_{n-i} q_1 \cdots q_{n-i-1}, \quad i = 0, 1, 2.
\]
We first show that
\begin{equation}
\label{eq:in1}
w_n < w_{n-2}.
\end{equation}

By simple rearrangement it follows that
\begin{align*}\label{e:condition_on_last_weights}
&
w_n<w_{n-2}\Longleftrightarrow
    q_{n-2}(1+q_{n-1}(1-q_n))<1.
\end{align*}

Using assumption~\eqref{eq:con1}, we immediately obtain
\begin{align*}
q_{n-2}\left(1 + q_{n-1}(1 - q_n)\right)
&< \frac{1}{\sqrt{2}} \left(1 + \frac{1}{\sqrt{2}}\left(1 - \left(1 - \frac{3 - \sqrt{5}}{2}\right)\right)\right)=0.89\cdots
\end{align*}
thus proving~\eqref{eq:in1}.


Now, let us consider the construction process of the OAT. Due to inequality~\eqref{eq:in1}, we have that, during the execution of the Hu--Tucker algorithm, the two nodes corresponding to weights \( w_n \) and \( w_{n-1} \) will be merged before \( w_{n-2} \) is involved in any merging operation. This is a key observation: it guarantees that there exists a step in the algorithm at which \( w_n \) and \( w_{n-1} \)---the two smallest and rightmost weights---are selected for merging.

As a consequence of this merging, both \( w_{n-1} \) and \( w_n \) will be assigned the same level \( \ell \geq 1 \) in the \emph{level assignment phase} of the Hu--Tucker algorithm. This level indicates their distance from the root in the final binary tree structure.

Since the nodes are ordered alphabetically (i.e., left to right from \( w_1 \) to \( w_n \)), the sequence of level assignments produced will be of the form
\[
\ell_1, \ell_2, \dots, \ell_{n-2}, \ell, \ell,
\]
where the last two equal entries \( \ell \) correspond to the merged weights \( w_{n-1} \) and \( w_n \). This level sequence is then used during the \emph{tree reconstruction phase}, in which the final structure of the OAT is obtained by recursively combining nodes according to their assigned levels, while preserving their left-to-right ordering.

Given that the two rightmost entries in the level sequence are equal, the reconstruction algorithm will combine the last two nodes into a common parent. This parent node will therefore appear as the right child of another internal node---specifically, it will be part of the \emph{right subtree of the root}. This follows from the recursive nature of reconstruction {(the rightmost sub-tree is formed from the highest-index elements in the sequence) and the fact that} OATs are always \emph{extended binary trees}\footnote{That is, every internal node, including the root, has exactly two children. {This guarantees that the root has the right subtree.}}. As a result, the \emph{initial test} (i.e., the root-level test in the associated decision process) must distinguish between units located in the left and right subtrees. Because \( w_{n-1} \) and \( w_n \) are located in the right subtree, the initial test must necessarily involve units corresponding to the left subtree---that is, among the units \( w_1, \dots, w_{n-2} \).

This completes the argument and establishes the claim.
\end{proof}

\end{document}
